\[ %%%%%%%%%%%%%%%%%%%%%%%%%%%%%%%%%%%%%%%%%%%%%%%%%%%%%%%%%%%%%%%%%%%%%%%%%%%%%%%% \section{Results and Analysis} \label{sec:results} %%%%%%%%%%%%%%%%%%%%%%%%%%%%%%%%%%%%%%%%%%%%%%%%%%%%%%%%%%%%%%%%%%%%%%%%%%%%%%%% This section presents the experimental evaluation of PHYSGEN and analyzes its performance across multiple nonlinear dynamical systems. The evaluation focuses on five central aspects: \begin{enumerate} \item long-horizon prediction accuracy; \item physical consistency of generated trajectories; \item stability during autonomous rollout; \item contribution of individual architectural components; \item generalization across different physical domains. \end{enumerate} The experiments are designed to validate the central hypothesis of PHYSGEN: \begin{quote} Explicit integration of physical constraints into graph-based dynamical learning improves long-horizon prediction accuracy, stability, and physical validity compared with existing data-driven and physics-informed approaches. \end{quote} %%%%%%%%%%%%%%%%%%%%%%%%%%%%%%%%%%%%%%%%%%%%%%%%%%%%%%%%%%%%%%%%%%%%%%%%%%%%%%%% \subsection{Experimental Overview} \label{sec:experimental_overview} %%%%%%%%%%%%%%%%%%%%%%%%%%%%%%%%%%%%%%%%%%%%%%%%%%%%%%%%%%%%%%%%%%%%%%%%%%%%%%%% The evaluation is performed using the benchmark suite described in Section~\ref{sec:benchmark_systems}: \begin{itemize} \item Lorenz-63 chaotic dynamical system; \item Reaction-Diffusion nonlinear PDE; \item Navier--Stokes fluid dynamics. \end{itemize} PHYSGEN is compared against representative baselines: \begin{itemize} \item MLP; \item LSTM; \item GNN; \item Graph Neural Simulator (GNS); \item Physics-Informed Neural Networks (PINNs); \item Fourier Neural Operator (FNO); \item DeepONet; \item Differentiable Physics models. \end{itemize} %%%%%%%%%%%%%%%%%%%%%%%%%%%%%%%%%%%%%%%%%%%%%%%%%%%%%%%%%%%%%%%%%%%%%%%%%%%%%%%% \subsubsection{Evaluation Protocol} %%%%%%%%%%%%%%%%%%%%%%%%%%%%%%%%%%%%%%%%%%%%%%%%%%%%%%%%%%%%%%%%%%%%%%%%%%%%%%%% Each model receives an initial physical state: \begin{equation} X_t \end{equation} and recursively predicts future states: \begin{equation} \hat{X}_{t+k}, \quad k=1,...,K. \end{equation} The evaluation is performed under identical conditions: \begin{itemize} \item identical training data; \item identical validation protocol; \item identical prediction horizon; \item identical evaluation metrics. \end{itemize} %%%%%%%%%%%%%%%%%%%%%%%%%%%%%%%%%%%%%%%%%%%%%%%%%%%%%%%%%%%%%%%%%%%%%%%%%%%%%%%% \subsubsection{Main Evaluation Dimensions} %%%%%%%%%%%%%%%%%%%%%%%%%%%%%%%%%%%%%%%%%%%%%%%%%%%%%%%%%%%%%%%%%%%%%%%%%%%%%%%% The results are analyzed according to four major dimensions. \paragraph{Prediction Accuracy.} The first dimension evaluates whether PHYSGEN improves forecasting accuracy: \begin{equation} RMSE = \sqrt{ \frac{1}{N} \sum_i (X_i-\hat{X}_i)^2 }. \end{equation} \paragraph{Physical Validity.} The second dimension evaluates whether predicted trajectories satisfy physical constraints: \begin{equation} \mathcal{F}(\hat{X})\approx0. \end{equation} \paragraph{Long-Horizon Stability.} The third dimension measures error growth during autonomous rollout: \begin{equation} E(k) = \| X_{t+k} - \hat{X}_{t+k} \|. \end{equation} \paragraph{Generalization Capability.} The fourth dimension evaluates transfer between different physical systems: \begin{equation} \mathcal{D}_i \rightarrow \mathcal{D}_j. \end{equation} %%%%%%%%%%%%%%%%%%%%%%%%%%%%%%%%%%%%%%%%%%%%%%%%%%%%%%%%%%%%%%%%%%%%%%%%%%%%%%%% \subsubsection{Expected Evaluation Trends} %%%%%%%%%%%%%%%%%%%%%%%%%%%%%%%%%%%%%%%%%%%%%%%%%%%%%%%%%%%%%%%%%%%%%%%%%%%%%%%% Based on the PHYSGEN design principles, we expect the following trends: \begin{itemize} \item purely data-driven models achieve competitive short-term prediction but experience rapid error accumulation; \item graph-based models improve representation of spatial interactions but lack explicit physical regulation; \item physics-informed approaches preserve constraints but may suffer from limited scalability; \item neural operators provide strong PDE approximation but do not explicitly control latent dynamical stability; \item PHYSGEN combines these advantages through physics-guided graph evolution and stability regularization. \end{itemize} %%%%%%%%%%%%%%%%%%%%%%%%%%%%%%%%%%%%%%%%%%%%%%%%%%%%%%%%%%%%%%%%%%%%%%%%%%%%%%%% \subsubsection{Statistical Evaluation} %%%%%%%%%%%%%%%%%%%%%%%%%%%%%%%%%%%%%%%%%%%%%%%%%%%%%%%%%%%%%%%%%%%%%%%%%%%%%%%% For each benchmark, experiments are repeated multiple times with different initial conditions. The reported performance is: \begin{equation} \bar{M} = \frac{1}{N} \sum_{i=1}^{N} M_i, \end{equation} with uncertainty estimated using: \begin{equation} \sigma_M = \sqrt{ \frac{1}{N} \sum_i (M_i-\bar{M})^2 }. \end{equation} %%%%%%%%%%%%%%%%%%%%%%%%%%%%%%%%%%%%%%%%%%%%%%%%%%%%%%%%%%%%%%%%%%%%%%%%%%%%%%%% \subsubsection{Summary} %%%%%%%%%%%%%%%%%%%%%%%%%%%%%%%%%%%%%%%%%%%%%%%%%%%%%%%%%%%%%%%%%%%%%%%%%%%%%%%% The experimental analysis evaluates PHYSGEN as a general framework for physics-guided long-horizon dynamical prediction. The following sections provide detailed comparisons of prediction accuracy, physical consistency, stability, ablation performance, and cross-domain generalization. \] 
